% ---------------------------------------------------------------------------
% Figure blocks for Paper IV.
%
% Generated by make_paper4_figures.py. Place each block in the section noted
% above it; main.tex already loads graphicx and sets \graphicspath, so the
% files can sit in the working directory or in ./figures/.
%
% The .pdf versions are vector and preferred for LaTeX; .png versions are
% provided for drafts and for journals that require raster.
% ---------------------------------------------------------------------------


% === place in Section 3 (Validation), after the Ryugu tables ===============
\begin{figure}[tp]
\centering
\includegraphics[width=\textwidth]{fig1.pdf}
\caption{Per-facet validation of the computed slope field against the
gravitational products released with the Hayabusa2 shape analysis of Ryugu,
on the identical $49\,152$-facet model at $\dens = 1200\ \mathrm{kg\,m^{-3}}$
and $P = 7.63$~h. (a) Facet-by-facet comparison; the colour scale is
logarithmic in the number of facets per bin, and the dashed line is the $1{:}1$
relation. (b) Distribution of the per-facet difference. The agreement is
scatter rather than offset: removing the mean bias leaves the RMS difference
unchanged.}
\label{fig:ryuguval}
\end{figure}


% === place in Section 3 (Validation), after the Bennu tables ===============
\begin{figure}[tp]
\centering
\includegraphics[width=\textwidth]{fig2.pdf}
\caption{Bennu. (a) Area-weighted mean and median slope by latitude band at the
measured bulk density; shaded bands mark the equatorial and high-latitude values
reported from OSIRIS-REx analyses. The profile reproduces the observed pattern
of low equatorial slopes rising to a mid-latitude maximum. (b) Dependence of the
slope on the assumed bulk density, resolved by latitude zone. Stars mark the
published global means at $0.85$ and $1.15\ \mathrm{g\,cm^{-3}}$. The
equatorial zone responds most strongly, as expected for a centrifugal effect
that scales with distance from the spin axis.}
\label{fig:bennuval}
\end{figure}


% === place in Section 5 (Geopotential dispersion), after Table \ref{tab:varseries} ===
\begin{figure}[tp]
\centering
\includegraphics[width=\textwidth]{fig3.pdf}
\caption{Behaviour of the two diagnostics under mesh refinement, expressed as
the fractional deviation from the value obtained on the finest mesh of each
sequence. (a) Area-weighted mean slope. (b) Normalized geopotential dispersion,
plotted on \emph{identical} axes. Both quantities are computed from the same
field on the same meshes, so the contrast is a property of the diagnostics
rather than of the calculation. The slope curves spread by up to a quarter of
their own value and are still moving at the finest meshes; the dispersion curves
are flat on this scale.}
\label{fig:resolution}
\end{figure}


% === place in Section 5, after Table \ref{tab:varsens} =====================
\begin{figure}[tp]
\centering
\includegraphics[width=\textwidth]{fig4.pdf}
\caption{Resolution diagnostics for the two quantities.
(a) Range sensitivity $S$, the peak-to-peak variation over the full mesh range
normalized by the mean. (b) Convergence residual $C$, the fractional change over
the final refinement step. The geopotential dispersion is the more stable diagnostic
on every body by both measures, but the two measures support claims of different
strength: the range advantage is modest on bodies whose slope is already stable,
whereas the convergence residual separates the two cleanly---the dispersion has
reached a plateau on every body, and the slope has not on Ryugu, Eros or
Itokawa.}
\label{fig:diagnostics}
\end{figure}


% === place in Section 6 (Cross-version comparison) =========================
\begin{figure}[tp]
\centering
\includegraphics[width=\textwidth]{fig5.pdf}
\caption{Shape-model detail and the slope distribution. (a) Slope distribution
of Itokawa from the spacecraft-derived ($N_f = 4\,857$) and radar-derived
($N_f = 3\,688$) shape models. The bulk descriptors, compared at a common facet
count by interpolation, agree closely---their area-weighted mean slopes differ
by $0.8\%$---but the radar model contains no steep tail.
(b) Effect of Laplacian smoothing at fixed facet count. Ten iterations bring the
spacecraft model of Itokawa to $12.90^{\circ}$, within $3.7\%$ of the value
obtained independently from the radar model (dashed line), supporting the
interpretation that the difference between the two models is dominated by
facet-scale relief.}
\label{fig:crossversion}
\end{figure}
